\begingroup
\small
\setlength{\tabcolsep}{3.5pt}
\renewcommand{\arraystretch}{1.15}
\begin{longtable}{@{}p{0.065\linewidth}p{0.285\linewidth}p{0.60\linewidth}@{}}
\caption{Hard-tier task catalogue. Method set A applies to every task in this table.}\label{tab:catalogue-hard}\\
\toprule
ID & Chemical task & Stated conditions and task-specific research target \\
\midrule
\endfirsthead
\multicolumn{3}{l}{\tablename\ \thetable\ (continued)}\\
\toprule
ID & Chemical task & Stated conditions and task-specific research target \\
\midrule
\endhead
\midrule
\multicolumn{3}{r}{Continued on the next page}\\
\endfoot
\bottomrule
\endlastfoot
HA1 & 3-methylcyclobutene ring opening: determine the stereochemical outcome\newline {\footnotesize Charge 0; multiplicity 1.} & \textit{Conditions:} gas phase or inert solvent, \(\sim\)150--200 \({}^{\circ}\mathrm{C}\), thermal.\newline \textit{Target:} Determine stereochemical products and channel identities; report channel barriers and the requested relative preference. \\
\addlinespace[3pt]
HA6 & trans-3,4-dimethylcyclobutene ring opening: determine the stereochemical outcome\newline {\footnotesize Charge 0; multiplicity 1.} & \textit{Conditions:} gas phase or inert solvent, heating, thermal.\newline \textit{Target:} Determine stereochemical products and channel identities; report channel barriers and the requested relative preference. \\
\addlinespace[3pt]
HA8 & (2E,4Z,6E)-octa-2,4,6-triene thermal 6-pi electrocyclization: determine the stereochemical outcome\newline {\footnotesize Charge 0; multiplicity 1.} & \textit{Conditions:} Gas phase; first thermal ground-state electrocyclization only; no subsequent hydrogen shifts or photochemistry.\newline \textit{Target:} Resolve the product stereochemistry of the first thermal electrocyclization; report investigated channel identities and barriers. \\
\addlinespace[3pt]
HB1 & cyclopentadiene + maleic anhydride Diels-Alder: determine endo/exo outcome\newline {\footnotesize Charge 0; multiplicity 1.} & \textit{Conditions:} benzene, room temperature or mild heating.\newline \textit{Target:} Determine endo/exo products, channel barriers, and relative preference. \\
\addlinespace[3pt]
HB10 & cyclopentadiene + acrylonitrile Diels-Alder reaction: determine both stereochemical pathways\newline {\footnotesize Charge 0; multiplicity 1.} & \textit{Conditions:} Gas phase; concerted addition; common reactant-complex reference for the two channels.\newline \textit{Target:} Compute both endo/exo pathways using a common reactant-complex reference; report identities and barriers. \\
\addlinespace[3pt]
HD2 & cyclohexene oxide + hydroxide: determine the stereochemical outcome of the epoxide opening\newline {\footnotesize Charge -1; multiplicity 1.} & \textit{Conditions:} aqueous NaOH, heat.\newline \textit{Target:} Determine the stereochemistry of the epoxide-opening product from the computed structures. \\
\addlinespace[3pt]
HF1 & cyclohex-2-en-1-ol peracid epoxidation: determine the stereochemical outcome\newline {\footnotesize Charge 0; multiplicity 1.} & \textit{Conditions:} performic acid, CH2Cl2, 0--25 \({}^{\circ}\mathrm{C}\).\newline \textit{Target:} Determine the two facial pathways, product stereochemistry, barriers, and relative preference. \\
\addlinespace[3pt]
HH1 & 1,1'-binaphthyl thermal racemization\newline {\footnotesize Charge 0; multiplicity 1.} & \textit{Conditions:} thermal racemization, \(\sim\)50 \({}^{\circ}\mathrm{C}\), DMF.\newline \textit{Target:} Characterize the thermal racemization pathway, stereochemical endpoints, and barrier. \\
\addlinespace[3pt]
HI1 & Cope elimination of threo/erythro N,N-dimethyl-3-phenyl-2-butylamine oxide: determine the stereochemical mapping\newline {\footnotesize Charge 0; multiplicity 1.} & \textit{Conditions:} Gas phase; isolated neutral singlet amine oxide. Treat threo and erythro substrates as two separate systems and retain explicit stereochemical SMILES.\newline \textit{Target:} Compute threo and erythro substrates separately; establish substrate CIP assignments and alkene E/Z identity from 3D geometries. Use each substrate's own minimum for its barrier. \\
\addlinespace[3pt]
HL1 & formonitrile oxide + norbornene [3+2] cycloaddition: determine both facial pathways\newline {\footnotesize Charge 0; multiplicity 1.} & \textit{Conditions:} Gas phase; neutral singlet formonitrile oxide and unsubstituted norbornene. Compare approach from both faces relative to the methano bridge.\newline \textit{Target:} Compute both facial approaches relative to the methano bridge and assign products from 3D geometries. Use each channel's own stable contact minimum; report contact-conformer energies separately. \\
\addlinespace[3pt]
\end{longtable}
\endgroup
